\documentclass[11pt, twocolumn, landscape, a4paper]{article}

\usepackage[margin=1cm]{geometry}
\usepackage[utf8]{inputenc}
\usepackage[spanish]{babel}
\usepackage{amssymb, amsmath, amsbsy}
\usepackage{tasks}
\usepackage{enumerate}
\usepackage{enumitem}
\usepackage{graphicx}

\usepackage[pdftex]{hyperref}
\hypersetup{pdftitle={Repaso en Polinomios / Valor numérico},pdfauthor={Luis Carrillo Gutiérrez}}

\fontfamily{cmss}\selectfont
\renewcommand{\familydefault}{\sfdefault}
\pagestyle{empty}

\begin{document}
\subsection*{Polinomios - V.N. (2)}

\begin{itemize}
\item{Si se define que : \quad$P(x - 3) = 4x - 7$\quad y \quad$P(F(x)) = 52x - 55$. \\ Calcular : $F(6)$
\begin{tasks}(5)
\task{56}\task{58}\task{60}\task{70}\task{72}
\end{tasks}
}
\item{Sabiendo que : \quad$P(x + 5) = 2x - 1$; \quad$P(Q(x) + 1) = 4x + 3$. \\ Calcular $Q(P(7))$
\begin{tasks}(5)
\task{10}\task{11}\task{12}\task{13}\task{14}
\end{tasks}
}
\item{Si se cumple que : \quad$G(F(x) - 2) = 3x + 1$, \quad además : \quad$G(x - 1) = x + 3$. \\ Hallar $G(F(2)) - 10$.
\begin{tasks}(5)
\task{-8}\task{-1}\task{1}\task{5}\task{9}
\end{tasks}
}
\item{Si : \quad$P(x + 3) = 5x + 7$, \quad además \quad$P[M(x) - 3] = 15x + 2$. \\ Calcular : $P[M(1)]$.
\begin{tasks}(5)
\task{35}\task{34}\task{33}\task{32}\task{31}
\end{tasks}
}
\item{Si se tiene que : \quad$P(x + 2) = 3x - 5$, \quad además \quad$P[H(x) - 3] = 9x + 16$. \\ Hallar $H(x + 5)$
\begin{tasks}(3)
\task{$3x - 11$}\task{$3x + 27$}\task{$3x - 1$}\task{$3x + 25$}\task{$27 - 3x$}
\end{tasks}
}
\end{itemize}
\end{document}